\section{Witnesses, Modes and Readers}
\label{sec:readers}

This section characterizes the countermodels of a disjunctive entailment over
arbitrary trees by the holes of finite automata. We first represent trees by
two sets of paths, in which the order becomes inclusion and a constraint
becomes a set of implications between atoms
(\cref{sec:readers-cones,sec:readers-atoms}). A failed query then has a
witness path, in one of two modes (\cref{sec:readers-modes}), and a family of
witnesses belongs to one solution exactly under three conditions on
derivations (\cref{sec:readers-realization}). Finite automata follow these
derivations along the witness paths (\cref{sec:readers-runs}), which gives
the readers and the reduction
(\cref{sec:readers-def,sec:readers-reduction}). The representation by leaf
cones is that of \cite{Zhao26E}; proofs omitted here are in
\cref{app:reduction}. Throughout, trees are trees over $\BFT$ and
assignments range over all trees.

\subsection{Leaf Cones}\label{sec:readers-cones}

\begin{definition}[Signs and cones]\label{def:cones}
The \emph{sign} of a leaf $v$ of a tree $t$ is $\sgp$ if $t(v)=\Top$ and
$\pol(v)=0$, or $t(v)=\Bot$ and $\pol(v)=1$; otherwise it is $\sgm$. For a
sign $s\in\{\sgp,\sgm\}$, the \emph{cone} of sign $s$ of $t$ is
\[
  \Cone st=\{\,w\in\Dir^*\mid\text{some prefix of $w$ is a leaf of $t$ of sign $s$}\,\}.
\]
We write $\flip s$ for the opposite sign.
\end{definition}

For a covariant constructor all polarities are $0$, so $\Cone{\sgp}t$ is the
set of paths along which $t$ reaches $\Top$ and $\Cone{\sgm}t$ the set of
paths along which it reaches $\Bot$; \cref{fig:tree-witness} shows an
example. With contravariant arguments the two labels exchange their roles at
paths of polarity $1$. A path has at most one prefix that is a leaf, and it
has none exactly when it is a position labeled $f$. A set $L$ of paths is
\emph{extension closed} if $w\in L$ implies $wi\in L$ for every $i\in\Dir$.

\begin{theorem}[Leaf representation]\label{thm:cones}
\begin{enumerate}
\item\label{it:cones-pair} For every tree $t$, the cones $\Cone{\sgp}t$ and
  $\Cone{\sgm}t$ are disjoint and extension closed. Conversely, every pair of
  disjoint extension-closed sets of paths is the pair of cones of exactly one
  tree.
\item\label{it:cones-order} For all trees $t$ and $u$:\quad
  $t\le u\iff\Cone{\sgp}t\subseteq\Cone{\sgp}u$ and
  $\Cone{\sgm}u\subseteq\Cone{\sgm}t$.
\item\label{it:cones-node} $\eps\in\Cone{\sgp}t$ iff $t=\Top$, and
  $\eps\in\Cone{\sgm}t$ iff $t=\Bot$. If $t=f(t_1,\dots,t_n)$, then for every
  $i\in\Dir$, path $w$ and sign $s$: $iw\in\Cone st$ iff
  $w\in\Cone{s}{t_i}$ when $i$ is covariant, and iff
  $w\in\Cone{\flip s}{t_i}$ when $i$ is contravariant.
\end{enumerate}
\end{theorem}

So going up in the order adds leaves of sign $\sgp$ and removes leaves of
sign $\sgm$. The variance enters only through the signs of the leaves; there
are two cones for each tree, whatever the arity.

%% tree-witness -- two trees and the witness path of a failed inequality.
\begin{figure}[t]
\centering
\begin{tikzpicture}[
  font=\small,
  nd/.style={inner sep=1.5pt, minimum size=4.5mm},
  mkd/.style={nd, draw, circle, inner sep=0.5pt},
  ed/.style={thin},
  wit/.style={line width=1.3pt},
  lb/.style={font=\scriptsize, inner sep=1pt}]
%% the tree of x
\begin{scope}
  \node[nd] (r) at (0,0) {$f$};
  \node[nd] (a) at (-0.9,-0.9) {$f$};
  \node[nd] (b) at (0.9,-0.9) {$\Bot$};
  \node[nd] (aa) at (-1.6,-1.8) {$\Bot$};
  \node[mkd] (ab) at (-0.2,-1.8) {$\Top$};
  \draw[wit] (r) -- node[lb, above left] {$1$} (a);
  \draw[ed] (r) -- node[lb, above right] {$2$} (b);
  \draw[ed] (a) -- node[lb, above left] {$1$} (aa);
  \draw[wit] (a) -- node[lb, above right] {$2$} (ab);
  \node at (0,0.55) {$\rho(x)$};
\end{scope}
%% the tree of y
\begin{scope}[xshift=4.3cm]
  \node[nd] (r) at (0,0) {$f$};
  \node[nd] (a) at (-0.9,-0.9) {$f$};
  \node[nd] (b) at (0.9,-0.9) {$\Top$};
  \node[nd] (aa) at (-1.6,-1.8) {$\Bot$};
  \node[mkd] (ab) at (-0.2,-1.8) {$f$};
  \node[nd] (aba) at (-0.7,-2.7) {$\Bot$};
  \node[nd] (abb) at (0.3,-2.7) {$\Bot$};
  \draw[wit] (r) -- node[lb, above left] {$1$} (a);
  \draw[ed] (r) -- node[lb, above right] {$2$} (b);
  \draw[ed] (a) -- node[lb, above left] {$1$} (aa);
  \draw[wit] (a) -- node[lb, above right] {$2$} (ab);
  \draw[ed] (ab) -- node[lb, above left] {$1$} (aba);
  \draw[ed] (ab) -- node[lb, above right] {$2$} (abb);
  \node at (0,0.55) {$\rho(y)$};
\end{scope}
\node[align=left, font=\footnotesize, anchor=west] at (6.3,-1.1)
  {$\Cone{\sgp}{\rho(x)}=12\,\Dir^*$\\[2pt]
   $\Cone{\sgm}{\rho(x)}=11\,\Dir^*\cup 2\,\Dir^*$\\[7pt]
   $\Cone{\sgp}{\rho(y)}=2\,\Dir^*$\\[2pt]
   $\Cone{\sgm}{\rho(y)}=11\,\Dir^*\cup121\,\Dir^*\cup122\,\Dir^*$};
\end{tikzpicture}
\caption{A witness of $\rho(x)\not\le\rho(y)$ for a binary covariant
constructor. Along the path $w=12$ (thick) the tree $\rho(x)$ reaches $\Top$
and $\rho(y)$ does not: $w\in\Cone{\sgp}{\rho(x)}\setminus\Cone{\sgp}{\rho(y)}$.
In the terms of \cref{sec:readers-modes} the query $(x,y)$ fails in mode
$\Top$ at $w$, with the seed $(x^{\sgp},w)$ and the exclusion
$(y^{\sgp},w)$.}
\label{fig:tree-witness}
\end{figure}


\subsection{Solutions as Closed Sets of Atoms}\label{sec:readers-atoms}

Let $\Gamma$ be a constraint over $V$. A \emph{signed variable} is a pair
$x^s$ of a variable and a sign; $\SV$ is the set of signed variables and
$\flip{x^s}=x^{\flip s}$. An \emph{atom} is a pair $(p,w)$ of a signed
variable and a path, and $\flip{(p,w)}=(\flip p,w)$ is its \emph{opposite}.
An assignment $\rho$ determines the set of atoms
\[
  S_\rho=\{\,(x^s,w)\mid w\in\Cone s{\rho(x)}\,\}.
\]
For an argument position $i$ we write $x^{[i]}$ for $x^{\sgp}$ if $i$ is
covariant and for $x^{\sgm}$ if $i$ is contravariant. \Cref{tab:clauses}
associates with every literal of $\Gamma$ a set of \emph{moves}, implications
between two atoms, together with \emph{facts}, atoms that must hold, and
\emph{forbidden atoms}, which must not. The \emph{rules} of $\Gamma$ are
its moves and the \emph{extensions} $(p,w)\to(p,wi)$, for all atoms $(p,w)$
and all $i\in\Dir$.

\begin{table}[t]
\caption{The rule form of the flat literals, for all paths $w$ and all
$i\in\Dir$.}
\label{tab:clauses}
\centering\small
\begin{tabular}{@{}llll@{}}
\toprule
Literal & Moves & Facts & Forbidden atoms \\
\midrule
$x\le y$ & $(x^{\sgp},w)\to(y^{\sgp},w)$,\quad $(y^{\sgm},w)\to(x^{\sgm},w)$ & & \\[2pt]
$f(x_1,\dots,x_n)\le y$ & $(x_i^{[i]},w)\to(y^{\sgp},iw)$,\quad
   $(y^{\sgm},iw)\to(\flip{x_i^{[i]}},w)$ & & $(y^{\sgm},\eps)$ \\[2pt]
$y\le f(x_1,\dots,x_n)$ & $(y^{\sgp},iw)\to(x_i^{[i]},w)$,\quad
   $(\flip{x_i^{[i]}},w)\to(y^{\sgm},iw)$ & & $(y^{\sgp},\eps)$ \\[2pt]
$x=\Bot$ & & $(x^{\sgm},\eps)$ & \\
$x=\Top$ & & $(x^{\sgp},\eps)$ & \\
\bottomrule
\end{tabular}
\end{table}

A set $S$ of atoms is \emph{closed} if it is closed under the rules,
\emph{consistent} if it contains no atom together with its opposite, and a
\emph{solution set} of $\Gamma$ if it is closed and consistent, contains the
facts and contains no forbidden atom.

\begin{lemma}[Rule form]\label{lem:clause-form}
The map $\rho\mapsto S_\rho$ is a bijection from the solutions of $\Gamma$
onto the solution sets of $\Gamma$.
\end{lemma}

The proof translates each literal with \cref{thm:cones}: for instance,
$f(x_1,\dots,x_n)\le y$ says that the cone of sign $\sgp$ of the constructor
tree, which consists of the paths $iw$ with $w$ in a cone of $\rho(x_i)$, is
contained in that of $\rho(y)$, and that the cone of sign $\sgm$ of
$\rho(y)$ is contained in that of the constructor tree, which does not
contain $\eps$.

We write $a\der b$ if the atom $b$ is \emph{derivable} from the atom $a$ by
finitely many rules, and $a\mv b$ if it is derivable by moves only. The
\emph{least closure} $\LC$ of $\Gamma$ is the set of atoms derivable from
facts. It is closed and contained in every solution set, so $\Gamma$ is
satisfiable if and only if $\LC$ is consistent and contains no forbidden
atom, and then $\LC$ is the least solution set.

The moves read the path of an atom as a stack whose top is its first letter:
a move replaces the signed variable and pops or pushes one letter at the top,
or leaves the path unchanged. Extensions add letters at the bottom. Two
properties of the moves are used throughout.

\begin{lemma}[Duality and normal form]\label{lem:duality}
\begin{enumerate}
\item\label{it:dual} If $a\to b$ is a move, so is $\flip b\to\flip a$. Hence
  $a\mv b$ implies $\flip b\mv\flip a$.
\item\label{it:ext-first} $(p,u)\der(q,v)$ if and only if
  $(p,u\sigma)\mv(q,v)$ for some path $\sigma$.
\end{enumerate}
\end{lemma}

Two finite sets of signed variables record where the forbidden atoms and the
facts enter. Both kinds of atoms have the empty path, and we let
\begin{equation}\label{eq:neg-fct}
\begin{aligned}
  \NegV&=\{\,r\in\SV\mid(r,\eps)\text{ is a forbidden atom of }\Gamma\,\},\\
  \FctV&=\{\,r\in\SV\mid(p,\eps)\mv(r,\eps)\text{ for some fact }(p,\eps)
  \text{ of }\Gamma\,\}.
\end{aligned}
\end{equation}
An atom whose derivations reach a signed variable of $\NegV$ with the empty
path cannot belong to a solution set, and an atom that is reached from a
signed variable of $\FctV$ with the empty path belongs to every solution
set.

\begin{example}[Rules of a small constraint]\label{ex:guard-rules}
Let $f$ be binary and covariant, and let
\[
  \Gamma=\{\,u\le f(u,u),\quad x\le u,\quad y=\Bot\,\}
\]
over the variables $x,y,u$. The inequality $u\le f(u,u)$ ensures that
$\Top$ does not occur in $u$ (this is the guard of
\cref{sec:decision-bf}); since $x\le u$, $\Top$ does not occur in $x$
either, and $y$ is $\Bot$. By \cref{tab:clauses} the moves are, for all paths $w$ and
$i\in\{1,2\}$,
\begin{align*}
  (u^{\sgp},iw)&\to(u^{\sgp},w), & (u^{\sgm},w)&\to(u^{\sgm},iw)
    &&\text{from }u\le f(u,u),\\
  (x^{\sgp},w)&\to(u^{\sgp},w), & (u^{\sgm},w)&\to(x^{\sgm},w)
    &&\text{from }x\le u.
\end{align*}
The first pops a letter, the second pushes one, and the last two keep the
path; in each row the second move is the dual of the first. The only forbidden atom is $(u^{\sgp},\eps)$ and the only fact is
$(y^{\sgm},\eps)$, so $\NegV=\{u^{\sgp}\}$ and $\FctV=\{y^{\sgm}\}$. The
least closure consists of the atoms $(y^{\sgm},w)$, which are derived from
the fact by extensions: it describes the solution with the least set of
atoms, $y=\Bot$ and $x=u=\fall$. The derivation
\[
  (x^{\sgp},12)\to(u^{\sgp},12)\to(u^{\sgp},2)\to(u^{\sgp},\eps)
\]
pops the path $12$ and ends in a forbidden atom. So no solution set contains
$(x^{\sgp},12)$: the tree of $x$ cannot reach $\Top$ along the path $12$,
nor along any other path, because every path can be popped in this way.
\end{example}

\subsection{Witnesses and Modes}\label{sec:readers-modes}

By \cref{thm:cones}\eqref{it:cones-order}, an inequality fails in one of two
ways.

\begin{lemma}[Witnesses]\label{lem:witness}
Let $\rho$ be an assignment and $(x,y)$ a query. Then $\rho(x)\not\le\rho(y)$
if and only if there are a path $w$ and a \emph{mode} $\mu\in\{\Top,\Bot\}$
such that
\begin{itemize}
\item $\mu=\Top$,\quad $(x^{\sgp},w)\in S_\rho$ and $(y^{\sgp},w)\notin
  S_\rho$;\quad or
\item $\mu=\Bot$,\quad $(y^{\sgm},w)\in S_\rho$ and $(x^{\sgm},w)\notin
  S_\rho$.
\end{itemize}
\end{lemma}

In mode $\Top$ the smaller side is too large along $w$: it has a leaf of
sign $\sgp$ that the larger side lacks. In mode $\Bot$ the larger side is too
small. We call the atom that must belong to $S_\rho$ the \emph{seed} of the
witness and the atom that must not belong to it the \emph{exclusion}; both
have the path $w$. For queries $(x_1,y_1),\dots,(x_k,y_k)$, a vector of modes
$\bar\mu=(\mu_1,\dots,\mu_k)$ and a tuple of paths $W=(w_1,\dots,w_k)$, we
write $A_i=(p_i,w_i)$ for the seed and $B_i=(p'_i,w_i)$ for the exclusion
of query $i$, so that
\begin{equation}\label{eq:seeds}
  (p_i,p'_i)=(x_i^{\sgp},y_i^{\sgp})\ \text{ if }\mu_i=\Top,\qquad
  (p_i,p'_i)=(y_i^{\sgm},x_i^{\sgm})\ \text{ if }\mu_i=\Bot.
\end{equation}
A solution $\rho$ of $\Gamma$ is a countermodel if and only if, for some
$\bar\mu$ and $W$, the set $S_\rho$ contains every $A_i$ and no $B_i$. The
signed variables $p_i,p'_i$ depend only on the queries and the modes; the
tuple $W$ is what the decision procedure has to find. In
\cref{fig:tree-witness}, for instance, the query $(x,y)$ fails in mode
$\Top$ with the witness $12$.

\subsection{Realization}\label{sec:readers-realization}

\begin{theorem}[Realization]\label{thm:realization}
Let $\Gamma$ be satisfiable over $\Trees$, and let $A_1,\dots,A_k$ and
$B_1,\dots,B_k$ be atoms. Some solution set of $\Gamma$ contains every $A_i$
and no $B_i$ if and only if
\begin{enumerate}[label=\textup{(R\arabic*)}]
\item\label{it:R1} no $A_i$ derives a forbidden atom;
\item\label{it:R2} $B_i\notin\LC$ and $\flip{A_i}\notin\LC$, for every $i$;
\item\label{it:R3} $A_j\not\der B_i$ and $A_j\not\der\flip{A_i}$, for all
  $i$ and $j$, including $i=j$.
\end{enumerate}
In this case the \emph{hull} $\LC\cup\{b\mid A_j\der b\text{ for some }j\}$
is the least such solution set.
\end{theorem}

The conditions are clearly necessary: a solution set is closed, so it
contains the hull, and it is consistent and contains no forbidden atom. For
the converse one has to show that the hull is consistent. If it contained an
atom $b$ derived from $A_j$ together with $\flip b$ derived from $A_i$, the
duality of \cref{lem:duality} would turn the second derivation around and
compose it with the first into a derivation from $A_j$ toward
$\flip{A_i}$. The extensions, which have no dual, are handled by
\ref{it:R1}.

Condition \ref{it:R3} is where the $k$ witnesses interact. It involves pairs
of queries, and also a single query with itself: a seed must not derive its
own exclusion or its own opposite.

\subsection{Runs along a Path}\label{sec:readers-runs}

The conditions of \cref{thm:realization} are conditions on derivations
between atoms whose paths are the unknown words $w_i$. Two families of
finite sets summarize the derivations by moves along a path.

\begin{definition}[Upper and lower runs]\label{def:runs}
For a signed variable $p$ and a path $w$,
\[
  \Up(p,w)=\{\,r\in\SV\mid(p,w)\mv(r,\eps)\,\},\qquad
  \Dn(p,w)=\{\,r\in\SV\mid(r,\eps)\mv(p,w)\,\}.
\]
\end{definition}

The \emph{upper run} $\Up(p,w)$ collects the signed variables reached from
$(p,w)$ when the whole path has been popped, and the \emph{lower run}
$\Dn(p,w)$ those from which $(p,w)$ is built by pushing the whole path. Since
moves only see the top of the stack, $r\in\Up(p,w)$ also means that
$(p,w\sigma)\mv(r,\sigma)$ for every $\sigma$ by moves that never pop a
letter of $\sigma$, and similarly for $\Dn$. By duality,
$\Dn(p,w)=\{\flip r\mid r\in\Up(\flip p,w)\}$.

\begin{lemma}[Runs are deterministic]\label{lem:runs-det}
There are computable functions $\stepU,\stepD\colon 2^{\SV}\times\Dir\to
2^{\SV}$ such that $\Up(p,wi)=\stepU(\Up(p,w),i)$ and
$\Dn(p,wi)=\stepD(\Dn(p,w),i)$ for all $p$, $w$ and $i$. Moreover
$\stepU(P,i)=\varnothing$ whenever $P\cap\NegV=\varnothing$.
\end{lemma}

So both runs are computed by reading $w$ from the root, with two
deterministic automata whose states are sets of signed variables. The last
statement holds because a letter can be popped only from an atom
$(y^{\sgp},iw)$ with a literal $y\le f(\dots)$, or from $(y^{\sgm},iw)$ with a
literal $f(\dots)\le y$, and in both cases the signed variable lies in
$\NegV$.

\begin{lemma}[Derivations through cuts]\label{lem:cuts}
Let $(p,u)$ and $(q,v)$ be atoms with $\Up(p,u)\cap\NegV=\varnothing$. Then
$(p,u)\der(q,v)$ if and only if there are cuts $a$ of $u$ and $b$ of $v$
such that
\[
  \wsuf ua\pre\wsuf vb\qquad\text{and}\qquad
  \Up(p,\wpre ua)\cap\Dn(q,\wpre vb)\ne\varnothing.
\]
Moreover, $(p,u)$ derives a forbidden atom if and only if
$\Up(p,u)\cap\NegV\ne\varnothing$.
\end{lemma}

A derivation from $(p,u)$ first extends $u$, then pops a prefix of the path
down to a lowest point and pushes a prefix of $v$. The hypothesis says that
the pops do not go beyond $u$, so the lowest point is a cut $a$ of $u$, and
the rest $\wsuf ua$ is never touched: it survives as a prefix of the rest of
$v$. This is the origin of the comparisons ``the rest of one word is a prefix
of the rest of another word'' in the definition of holes. In
\cref{ex:guard-rules}, the atom $(u^{\sgm},2)$ derives $(x^{\sgm},12)$ by
pushing the letter $1$ and moving to $x^{\sgm}$. The cuts of the lemma are
$0$ in the path $2$ and $1$ in the path $12$: the rest $2$ is not touched,
and $\Up(u^{\sgm},\eps)=\{u^{\sgm},x^{\sgm}\}$ meets
$\Dn(x^{\sgm},1)=\{u^{\sgm}\}$. The same analysis,
applied to derivations from facts, shows that $(q,v)\in\LC$ if and only if
$\Dn(q,\wpre vb)$ meets the set $\FctV$ of \eqref{eq:neg-fct} for some cut
$b$ of $v$; this uses that $\Gamma$ is satisfiable.

\subsection{Structural Readers and Holes}\label{sec:readers-def}

We now define the automata and their holes abstractly. The alphabet is an
arbitrary finite set $\Alp$; in the reduction it is $\Dir$.

\begin{definition}[Structural reader]\label{def:reader}
A \emph{structural reader} with $k$ \emph{starts} over $\Alp$ consists of a
finite set $S$ of states, a transition function $\delta\colon S\times\Alp\to
S$, $k$ start states $s_1,\dots,s_k\in S$, which need not be distinct, a
\emph{dead}
state $\dead\in S$ with $\delta(\dead,c)=\dead$ for every letter $c$, a set
$\Rej\subseteq S\setminus\{\dead\}$ of \emph{rejected} states and a relation
$\Adm\subseteq S\times S$ of \emph{admitted} pairs, such that for every
$q\in\Rej$ and every letter $c$ with $\delta(q,c)\ne\dead$:
\begin{enumerate}[label=\textup{(S\arabic*)}]
\item\label{it:S1} $\delta(q,c)\in\Rej$;
\item\label{it:S2} $(q',\delta(q,c))\notin\Adm$ for every state $q'$.
\end{enumerate}
Its \emph{size} is $N=\len S$.
\end{definition}

The names anticipate the reduction of \cref{sec:readers-reduction}, where a
reader follows the derivations from the seed of a witness. A run is
\emph{dead} when the constraint alone already contradicts the witness. A
state is \emph{rejected} when the seed can no longer derive a forbidden
atom: the reader has rejected every way of refuting the witness by a
forbidden atom. So a rejected state is what a witness has to reach, although
the word may suggest the opposite. Condition \ref{it:S1} says that rejection
is permanent, and \ref{it:S2} that a state reached after a rejected state is
never the second component of an admitted pair. A pair of states is
\emph{admitted} when it shows a derivation from the seed of one witness
that contradicts a witness, the same one or another; such pairs matter only
at the places defined next.

Let $W=(w_1,\dots,w_k)$ be a tuple of words. The word $w_i$ is read from
$s_i$; we write $q_i(a)=\delta^*(s_i,\wpre{w_i}a)$ for the state of
coordinate $i$ at the cut $a$, and $q_i(\len{w_i})$ is its \emph{endpoint}.
The \emph{total length} of $W$ is $\tot W=\len{w_1}+\dots+\len{w_k}$.

\begin{definition}[Comparisons, channels and holes]\label{def:hole}
Let $i$ and $j$ be coordinates, possibly equal. A \emph{comparison} from the
\emph{host} $i$ to the \emph{target} $j$ is a pair $(a,b)$ of a cut $a$ of
$w_i$ and a cut $b$ of $w_j$ such that $\wsuf{w_j}b\pre\wsuf{w_i}a$: the
rest of the target is a prefix of the rest of the host. The pair of states
$(q_i(a),q_j(b))$ is the \emph{channel pair} at this comparison, and the
\emph{channel} from $i$ to $j$ is the set of all channel pairs,
\[
  \Chan_{ij}(W)=\{\,(q_i(a),q_j(b))\mid a\text{ a cut of }w_i,\ b\text{ a cut
  of }w_j,\ \wsuf{w_j}b\pre\wsuf{w_i}a\,\}.
\]
Coordinate $i$ is \emph{eligible} if $q_i(a)\ne\dead$ for every cut $a$ and
its endpoint is rejected. The tuple $W$ is a \emph{hole} if every coordinate
is eligible and $\Chan_{ij}(W)\cap\Adm=\varnothing$ for all $i$ and $j$.
\end{definition}

All cuts are included, also the two ends: at $b=\len{w_j}$ the rest of the
target is empty, so $(a,b)$ is a comparison for every cut $a$ of the host.
For example, in the tuple $(0110,\,10)$ the comparisons from the host $1$ to
the target $2$ are $(2,0)$, since $10$ is the rest of $0110$ after the cut
$2$, then $(0,1)$ and $(3,1)$, since the rest $0$ of the target is a prefix
of the rests $0110$ and $0$, and $(a,2)$ for the five cuts $a$ of the host.
In a pair $(q,q')\in\Adm$ the first component is the state of a host and the
second the state of a target; the relation need not be symmetric.
\Cref{fig:reader-hole}(a) shows a comparison and its channel pair.

For an eligible coordinate, \ref{it:S1} gives a \emph{first rejected cut}
$J_i$: the states $q_i(a)$ are rejected exactly for $a\ge J_i$. The prefix
$\wpre{w_i}{J_i}$ is the \emph{head} of $w_i$ and the suffix
$\wsuf{w_i}{J_i}$ its \emph{rejected tail}. By \ref{it:S2}, the target cut of
an admitted channel pair is at most $J_j$. So the tail of a word is never
involved as a target, but it can be involved as a host, and for $i=j$ a
channel pair has $a\le b$.

%% reader-hole -- a channel between two words of a tuple, and the reader of the worked example.
\begin{figure}[t]
\centering
\begin{tikzpicture}[font=\small, >={Stealth[length=2mm]}]
%% Left: a channel between a host word and a target word.
\begin{scope}
  \def\h{0.36}
  % host word
  \draw (0,0) rectangle (6.4,\h);
  \fill[black!12] (2.2,0) rectangle (5.0,\h);
  \draw (2.2,0) -- (2.2,\h);
  \draw (5.0,0) -- (5.0,\h);
  \node[left] at (0,0.5*\h) {$w_i$};
  \node[above, font=\scriptsize] at (2.2,\h) {cut $a$};
  \node[font=\scriptsize] at (3.6,0.5*\h) {$\wsub{w_i}{a}{a+\len{w_j}-b}$};
  % target word
  \draw (0.6,-1.5) rectangle (5.0,-1.5+\h);
  \fill[black!12] (2.2,-1.5) rectangle (5.0,-1.5+\h);
  \draw (2.2,-1.5) -- (2.2,-1.5+\h);
  \node[left] at (0.6,-1.5+0.5*\h) {$w_j$};
  \node[below, font=\scriptsize] at (2.2,-1.5) {cut $b$};
  \node[font=\scriptsize] at (3.6,-1.5+0.5*\h) {$\wsuf{w_j}{b}$};
  \draw[densely dotted] (2.2,0) -- (2.2,-1.5+\h);
  \draw[densely dotted] (5.0,0) -- (5.0,-1.5+\h);
  \node[font=\scriptsize] at (3.6,-0.57) {equal};
  % states
  \node[circle, draw, inner sep=1pt, font=\scriptsize] (qa) at (2.2,1.25) {$q_i(a)$};
  \node[circle, draw, inner sep=1pt, font=\scriptsize] (qb) at (2.2,-2.7) {$q_j(b)$};
  \draw[->] (qa) -- (2.2,\h+0.32);
  \draw[->] (qb) -- (2.2,-1.5-0.32);
  \node[font=\scriptsize, align=left, anchor=west] at (3.1,1.25)
    {$(q_i(a),q_j(b))\in\Chan_{ij}(W)$};
  \node[font=\scriptsize, align=center] at (3.2,-3.45) {(a) a comparison $(a,b)$: host $i$, target $j$};
\end{scope}
%% Right: the reader of the worked example.
\begin{scope}[xshift=8.0cm, yshift=-0.2cm,
    st/.style={draw, circle, double, double distance=1pt, minimum size=7mm, inner sep=0pt, font=\scriptsize}]
  \node[st] (s1) at (0,0.7) {$s_1$};
  \node[st] (s2) at (0,-1.6) {$s_2$};
  \node[st] (z) at (2.6,-0.45) {$z$};
  \draw[<-] (s1) -- ++(-0.9,0);
  \draw[<-] (s2) -- ++(-0.9,0);
  \draw[->] (s1) -- node[above, sloped, font=\scriptsize] {$1,\dots,n$} (z);
  \draw[->] (s2) -- node[below, sloped, font=\scriptsize] {$1,\dots,n$} (z);
  \draw[->] (z) edge[loop above] node[font=\scriptsize] {$1,\dots,n$} (z);
  \draw[dashed, <->] (s1) -- node[left, font=\scriptsize] {$\Adm$} (s2);
  \node[font=\scriptsize, align=left, anchor=west] at (-1.3,2.1)
    {$s_1=(\{x^{\sgp}\},\{y^{\sgp}\})$,\quad $s_2=(\{x^{\sgm}\},\{y^{\sgm}\})$,\\[1pt] $z=(\varnothing,\varnothing)$};
  \node[font=\scriptsize, align=center] at (1.3,-3.25) {(b) the reader of \cref{ex:reader-unary}};
\end{scope}
\end{tikzpicture}
\caption{(a) A comparison of a tuple $W$ and its channel pair: the rest of
the target word $w_j$ after the cut $b$ is a prefix of the rest of the host
word $w_i$ after the cut $a$, and the pair of the states at the two cuts
belongs to the channel from $i$ to $j$. The tuple is a hole only if no such
pair is admitted, for all $i$ and $j$, including $i=j$. (b) The reader of
the empty constraint with the queries $(x,y)$ and $(y,x)$ in the mode vector
$(\Top,\Bot)$. All three states are rejected (double circles), the dead
state is not shown, and the admitted pairs are $(s_1,s_2)$ and
$(s_2,s_1)$.}
\label{fig:reader-hole}
\end{figure}


\subsection{The Reader of a Mode and the Reduction}\label{sec:readers-reduction}

Fix a satisfiable constraint $\Gamma$, queries $(x_1,y_1),\dots,(x_k,y_k)$
and a vector of modes $\bar\mu$, with the signed variables $p_i,p'_i$ of
\eqref{eq:seeds} and the sets $\NegV$ and $\FctV$ of \eqref{eq:neg-fct}.

\begin{definition}[The reader $M_{\bar\mu}$]\label{def:compile}
The states of $M_{\bar\mu}$ are the pairs $(P,Q)$ of sets of signed variables
with $(Q\cup\flip P)\cap\FctV=\varnothing$, where $\flip P=\{\flip r\mid r\in
P\}$, and a dead state $\dead$. The transition on a letter $i$ takes $(P,Q)$
to $(\stepU(P,i),\stepD(Q,i))$ if this pair is a state, and to $\dead$
otherwise. The $i$th start state is $(\Up(p_i,\eps),\Dn(p'_i,\eps))$ if this
pair is a state, and $\dead$ otherwise. Moreover
\begin{align*}
  (P,Q)\in\Rej&\iff P\cap\NegV=\varnothing,\\
  ((P_h,Q_h),(P_t,Q_t))\in\Adm&\iff P_t\cap(Q_h\cup\flip{P_h})\ne\varnothing.
\end{align*}
\end{definition}

Coordinate $i$ follows the upper run of its seed, in the \emph{upper
component} $P$, and the lower run of its exclusion, in the \emph{lower
component} $Q$. It dies when a prefix of
the path shows that the exclusion or the opposite of the seed lies in the
least closure, so that every solution contradicts the witness. It is
rejected when its seed derives no forbidden atom. A pair is admitted when
the seed of the target derives the exclusion of the host or the opposite of
the seed of the host, at a comparison of the two paths as in
\cref{lem:cuts}. With these readings the three conditions of
\cref{thm:realization} are exactly the conditions of a hole: \ref{it:R2} says
that no run dies, \ref{it:R1} that every endpoint is rejected, and
\ref{it:R3} that no channel contains an admitted pair.

%% guard-reader -- the two readers of the worked example with a nonempty constraint.
\begin{figure}[t]
\centering
\begin{tikzpicture}[font=\scriptsize, >={Stealth[length=2mm]},
    st/.style={draw, circle, minimum size=7mm, inner sep=0pt},
    rj/.style={st, double, double distance=1pt}]
%% mode top
\begin{scope}
  \node[st] (t0) at (0,0) {$t_0$};
  \node[st] (t1) at (2.2,0) {$t_1$};
  \draw[<-] (t0) -- ++(-0.9,0);
  \draw[->] (t0) -- node[above] {$1,2$} (t1);
  \draw[->] (t1) edge[loop right] node {$1,2$} (t1);
  \node[anchor=west, align=left] at (-0.9,-1.05)
    {$t_0=(\{x^{\sgp},u^{\sgp}\},\{y^{\sgp}\})$\\[1pt] $t_1=(\{u^{\sgp}\},\varnothing)$};
  \node at (1.3,-2.05) {(a) mode $\Top$: no state is rejected};
\end{scope}
%% mode bot
\begin{scope}[xshift=7.0cm]
  \node[rj] (s0) at (0,0) {$s_0$};
  \node[rj] (s1) at (2.2,0) {$s_1$};
  \draw[<-] (s0) -- ++(-0.9,0);
  \draw[->] (s0) -- node[above] {$1,2$} (s1);
  \draw[->] (s1) edge[loop right] node {$1,2$} (s1);
  \node[anchor=west, align=left] at (-0.9,-1.05)
    {$s_0=(\{y^{\sgm}\},\{x^{\sgm},u^{\sgm}\})$\\[1pt] $s_1=(\varnothing,\{u^{\sgm}\})$};
  \node at (1.3,-2.05) {(b) mode $\Bot$: both states are rejected};
\end{scope}
\end{tikzpicture}
\caption{The readers of $\Gamma=\{u\le f(u,u),\ x\le u,\ y=\Bot\}$ and the
query $(x,y)$ (\cref{ex:guard-reader}). A state is the pair of the upper run
of the seed and the lower run of the exclusion; rejected states are drawn
with a double circle, and the dead state, which is not reachable here, is
not shown. (a) In mode $\Top$ the upper component always contains
$u^{\sgp}\in\NegV$, so there is no hole. (b) In mode $\Bot$ every word is a
hole, since no pair of states is admitted.}
\label{fig:guard-reader}
\end{figure}


\begin{theorem}[Reduction]\label{thm:reduction}
$M_{\bar\mu}$ is a structural reader over $\Dir$ with $k$ starts and at most
$2^{4\len V}+1$ states, computable from $\Gamma$, the queries and $\bar\mu$.
A tuple $W$ of paths is a hole of $M_{\bar\mu}$ if and only if some solution
of $\Gamma$ contains the seeds and none of the exclusions of $(\bar\mu,W)$.
Consequently, for a constraint $\Gamma$ that is satisfiable over $\Trees$,
\[
  \Gamma\not\entails_{\Trees}x_1\le y_1\vee\dots\vee x_k\le y_k
  \iff
  \text{$M_{\bar\mu}$ has a hole for some $\bar\mu\in\{\Top,\Bot\}^k$.}
\]
\end{theorem}

\begin{proof}[Proof sketch]
By \cref{lem:runs-det}, the state of an eligible coordinate $i$ at a cut $a$
is the pair $(\Up(p_i,\wpre{w_i}a),\Dn(p'_i,\wpre{w_i}a))$. Once $P\cap\NegV$ is empty,
the next upper set is empty, which gives \ref{it:S1} and \ref{it:S2}. A run
stays alive exactly when no prefix puts $B_i$ or $\flip{A_i}$ into the least
closure, which is \ref{it:R2}, and its endpoint is rejected exactly when
\ref{it:R1} holds. Under \ref{it:R1}, \cref{lem:cuts} states that $A_j\der
B_i$ or $A_j\der\flip{A_i}$ if and only if there are cuts $b$ of $w_j$ and
$a$ of $w_i$ with $\wsuf{w_j}b\pre\wsuf{w_i}a$ and
$\Up(p_j,\wpre{w_j}b)\cap(\Dn(p'_i,\wpre{w_i}a)\cup\Dn(\flip{p_i},\wpre{w_i}a))
\ne\varnothing$, that is, an admitted pair on the channel from $i$ to $j$.
The last statement follows with \cref{lem:witness}.
\end{proof}

\begin{example}[One query: the reader of a small constraint]\label{ex:guard-reader}
Consider the constraint $\Gamma=\{u\le f(u,u),\ x\le u,\ y=\Bot\}$ of
\cref{ex:guard-rules}, with $\NegV=\{u^{\sgp}\}$ and $\FctV=\{y^{\sgm}\}$,
and the query $(x,y)$. \Cref{fig:guard-reader} shows the two readers.

\emph{Mode $\Top$.} The seed is $(x^{\sgp},w)$ and the exclusion
$(y^{\sgp},w)$: the tree of $x$ should reach $\Top$ along $w$. The upper run
of the seed is $\Up(x^{\sgp},\eps)=\{x^{\sgp},u^{\sgp}\}$ and
$\Up(x^{\sgp},w)=\{u^{\sgp}\}$ for $w\ne\eps$, because the only way to pop
the path of the seed is to move to $u^{\sgp}$ and to pop there. No move ends
in $(y^{\sgp},w)$, so the lower run of the exclusion is $\{y^{\sgp}\}$ at
$\eps$ and empty afterwards. The upper component always contains
$u^{\sgp}\in\NegV$, so no state of the run is rejected, and the reader has
no hole: the seed always derives the forbidden atom $(u^{\sgp},\eps)$, as in
\cref{ex:guard-rules}.

\emph{Mode $\Bot$.} The seed is $(y^{\sgm},w)$ and the exclusion
$(x^{\sgm},w)$: the tree of $y$ should reach $\Bot$ along $w$ and the tree
of $x$ should not. No move starts from $y^{\sgm}$, so the upper run is
$\{y^{\sgm}\}$ at $\eps$ and empty afterwards. The lower run of the
exclusion is $\Dn(x^{\sgm},\eps)=\{x^{\sgm},u^{\sgm}\}$ and
$\Dn(x^{\sgm},w)=\{u^{\sgm}\}$ for $w\ne\eps$, since $(x^{\sgm},w)$ is
built from $(u^{\sgm},\eps)$ by pushing $w$ at $u^{\sgm}$ and moving to
$x^{\sgm}$. Both states are rejected, neither is dead, because
$y^{\sgm}\in\FctV$ belongs neither to the lower components nor to the
flipped upper components, and no pair is admitted: for the only nonempty
upper component, $\{y^{\sgm}\}\cap(\{x^{\sgm},u^{\sgm}\}\cup\{y^{\sgp}\})=
\varnothing$. So every word is a hole, in particular $w=\eps$. The hull of
its seed $(y^{\sgm},\eps)$ is the least closure, the solution $y=\Bot$,
$x=u=\fall$, and indeed $\fall\not\le\Bot$: the entailment
$\Gamma\entails x\le y$ fails.

\emph{A dead run.} For the reversed query $(y,x)$, the start state of mode
$\Top$ would be $(\{y^{\sgp}\},\{x^{\sgp}\})$, whose flipped upper component
contains the fact $y^{\sgm}$; and the start state of mode $\Bot$ would be
$(\{x^{\sgm}\},\{y^{\sgm}\})$, whose lower component contains it. So both
readers start in $\dead$ and have no hole: $\Gamma$ entails $y\le x$, as it
must, since $y=\Bot$ in every solution.
\end{example}

\begin{example}[Two queries, unary and binary]\label{ex:reader-unary}
Take the empty constraint and the queries $(x,y)$ and $(y,x)$ of
\cref{ex:unary}. There are no moves, no facts and no forbidden atoms, so
$\Up(p,\eps)=\Dn(p,\eps)=\{p\}$, both sets are empty for nonempty paths,
nothing is dead and every state is rejected. In the mode vector $(\Top,\Bot)$
the seeds are $(x^{\sgp},w_1)$ and $(x^{\sgm},w_2)$ and the exclusions
$(y^{\sgp},w_1)$ and $(y^{\sgm},w_2)$: the tree of $x$ should have a leaf of
sign $\sgp$ on $w_1$ and a leaf of sign $\sgm$ on $w_2$. The reader is shown
in \cref{fig:reader-hole}(b). The pair $(s_1,s_2)$ is admitted because the
upper component $\{x^{\sgm}\}$ of the target $s_2$ meets the flipped upper
component $\flip{\{x^{\sgp}\}}=\{x^{\sgm}\}$ of the host $s_1$, and
symmetrically for $(s_2,s_1)$.

A channel pair $(s_1,s_2)$ arises from the comparison $(0,0)$, which exists
when $w_2\pre w_1$, and $(s_2,s_1)$ when $w_1\pre w_2$. So $(w_1,w_2)$ is a hole if and only if
$w_1$ and $w_2$ are incomparable: two leaves of opposite signs cannot lie on
one path. For $n=1$ any two words are comparable and there is no hole; the
other three mode vectors give the same reader up to renaming, so the
entailment holds. For $n=2$ the tuple $(1,2)$ is a hole. The hull of its
seeds is the solution $x=f(\Top,\Bot)$, $y=\fall$ for a covariant
constructor, a countermodel.
\end{example}
